\chapter{Implementación y reproducibilidad}
\label{ap:software}

\section{Código}

Toda la implementación es en Python 3, con \texttt{numpy}, \texttt{pandas},
\texttt{matplotlib} y \texttt{streamlit}. El código completo --- el motor
actuarial con las derivadas exactas, los \datoNMetodos{} métodos iterativos con su
instrumentación, el banco de pruebas, los generadores de cifras, tablas y figuras,
y la aplicación interactiva de la defensa --- forma parte de los archivos del
proyecto. Este identifica como repositorio del trabajo:
\begin{center}
\repositorio
\end{center}
La reproducción requiere la versión del código, los datos y la configuración
asociados a cada ejecución, junto con las dependencias del entorno. El enlace por
sí solo no identifica esa revisión ni acredita la disponibilidad pública de todas
las entradas y salidas. El registro de una ejecución debe conservar su fecha y la
revisión o las sumas de comprobación de los archivos utilizados.

% ---------------------------------------------------------------------
%  Detalle de modulos: trasladado al README del repositorio (31-oct-2026).
%  Para reincorporarlo, borrar las dos lineas que lo encierran (iffalse ... fi).
% ---------------------------------------------------------------------
\iffalse
\section{Arquitectura}

Toda la implementación es en Python 3, con \texttt{numpy} para el cálculo
vectorizado, \texttt{pandas} para el manejo de las bases de pólizas,
\texttt{matplotlib} para las figuras y \texttt{streamlit} para la aplicación
interactiva. El proyecto identifica el repositorio \repositorio; los módulos de
la implementación se organizan como sigue.

\begin{center}
\small
\begin{tabular}{@{}p{4.6cm}p{9.6cm}@{}}
\toprule
\textbf{Módulo} & \textbf{Contenido} \\
\midrule
\multicolumn{2}{@{}l}{\textit{Motor actuarial (preexistente, extendido en este trabajo)}}\\
\texttt{core/config.py} & parámetros actuariales y constantes normativas \\
\texttt{core/mortalidad.py} & series de mortalidad, mensualización y factores de mejora; fuentes por conciliar \\
\texttt{core/motor.py} & función de valoración $\VP$; \textbf{derivadas laterales exactas} (aporte de este trabajo) \\
\texttt{core/smith\_wilson.py} & extrapolación de la curva de descuento \\
\texttt{core/retiro\_programado.py} & recursión de retiro programado y operadores de ajuste \\
\texttt{core/base\_masiva.py} & carga y procesamiento de bases de pólizas \\
\midrule
\multicolumn{2}{@{}l}{\textit{Aporte de esta tesis}}\\
\texttt{core/metodos\_iterativos.py} & los \datoNMetodos{} métodos con interfaz común e instrumentación \\
\texttt{core/problema\_tesis.py} & construcción de $f$, $\dB f$ y de los escenarios \\
\texttt{core/banco\_pruebas.py} & ejecución del catálogo y agregación de resultados \\
\midrule
\multicolumn{2}{@{}l}{\textit{Generación del documento}}\\
\texttt{correr\_benchmark.py} & ejecuta el banco de pruebas completo \\
\texttt{fig\_01..05\_*.py} & genera las figuras en PDF vectorial \\
\texttt{generar\_macros\_tex.py} & genera \texttt{datos\_numericos.tex} \\
\texttt{generar\_tablas\_tex.py} & genera las tablas de resultados en \LaTeX \\
\texttt{generar\_extractos\_codigo.py} & extrae los listados del código fuente \\
\texttt{insertar\_notas.py} & inserta las notas al pie verificando cada ancla \\
\midrule
\multicolumn{2}{@{}l}{\textit{Aplicación interactiva}}\\
\texttt{app/} & presentación de la defensa en Streamlit \\
\bottomrule
\end{tabular}
\end{center}

\fi

\section{Sobre la trazabilidad de las cifras}

Los scripts
\texttt{generar\_macros\_tex.py} y \texttt{generar\_tablas\_tex.py} producen,
respectivamente, un fichero de macros \LaTeX{} y los cuerpos de las tablas, ambos a
partir de los resultados del banco de pruebas. El texto utiliza esas macros para
las magnitudes cubiertas por los generadores. Este mecanismo reduce la
transcripción manual, pero requiere volver a ejecutar el banco y los generadores
cuando cambian las entradas o el código. Las figuras, los listados y las
afirmaciones narrativas deben contrastarse también con esas salidas: la generación
no garantiza por sí sola la coherencia del documento. Las notas al pie y sus
anclas requieren revisión de contenido aunque se inserten mediante un script.

La contabilidad distingue la preparación y el bucle iterativo, y registra las
llamadas efectivas al motor. Cada ejecución de un método dispone de su propia
memoria de evaluaciones. Cuando se solicita una derivada analítica, un oráculo
conjunto obtiene reserva y derivadas en una misma llamada; las diferencias finitas
conservan el costo de sus valoraciones auxiliares. El tiempo medido complementa
estos contadores y depende del entorno de ejecución.

\section{Los tres componentes de la derivada exacta}

La contribución algorítmica principal se concentra en una extensión de
\texttt{core/motor.py}. Al invocar \texttt{calcular\_reserva} con
\texttt{calcular\_derivada=True}, la función devuelve además las dos derivadas
laterales de \cref{eq:subdif-der,eq:subdif-izq}. El cálculo reutiliza los vectores
construidos para evaluar $\VP$ y añade máscaras y productos escalares. Evita las
valoraciones auxiliares necesarias para estimar la pendiente por diferencias
finitas, pero tiene un costo marginal que debe medirse. Una llamada conjunta no
significa que la derivada sea gratuita ni que todo un método cueste una única
valoración por iteración, pues también hay preparación y comprobaciones.

Los tres ingredientes son:
\begin{enumerate}[leftmargin=*,nosep]
  \item la máscara de actividad del piso, $a_tM\gtrless b_t$, con desigualdad
        estricta o no según el lado;
  \item la máscara de banda del auxilio funerario,
        $5b_t\lessgtr a_tM\lessgtr10b_t$;
  \item la corrección de la rama de crecimiento en $M=S_0$
        (\cref{prop:continuidad-rama}), sin la cual la derivada en ese punto es
        incorrecta; el factor $34$ corresponde al ejemplo numérico del \cref{cap:estructura}. La corrección afecta a los métodos que consultan
        la derivada en ese punto; bisección y secante no la necesitan.
\end{enumerate}

% Comandos de reproduccion: trasladados al README del repositorio.
\iffalse
\section{Reproducción de los resultados}

\begin{verbatim}
# 1. Dependencias
pip install numpy pandas matplotlib openpyxl scipy streamlit plotly

# 2. Banco de pruebas completo
cd py && python3 correr_benchmark.py

# 3. Figuras
python3 fig_01_estructura.py
python3 fig_02_convergencia.py
python3 fig_03_actuarial.py
python3 fig_04_rp_sensibilidad.py
python3 fig_05_variables.py

# 4. Cifras, tablas y listados del documento
python3 generar_macros_tex.py
python3 generar_tablas_tex.py
python3 generar_extractos_codigo.py

# 5. Compilación
cd ../tex && latexmk -pdf tesis.tex

# 6. Aplicación interactiva de la defensa
cd ../app && streamlit run presentacion.py
\end{verbatim}

\fi

\section{Semillas y reproducibilidad}

La simulación del rendimiento de la cuenta individual admite una semilla
explícita. Cuando no se proporciona, la implementación deriva un entero de un
resumen SHA--256 de entradas serializadas en JSON, con orden y formato fijos.
Esta construcción evita la variación entre procesos del resumen interno de
Python. No garantiza semillas distintas para todas las entradas posibles ni
demuestra que asignar trayectorias diferentes a distintas pólizas sea una
representación adecuada del riesgo de mercado.

Reproducir una trayectoria requiere conservar semilla, generador, entradas,
configuración y versiones de las dependencias; las diferencias de entorno pueden
afectar los tiempos y el redondeo. Para comparar variantes de una póliza se debe
fijar una misma semilla o una misma trayectoria de rendimientos, pues una semilla
derivada de entradas que cambian puede introducir simultáneamente una nueva
realización de mercado. La reproducibilidad computacional no acredita la
calibración del proceso estocástico.
